% The body of the document. Each section demonstrates one feature once:
% one equation, one booktabs table, one Graphviz diagram and one code block.
% Replace every sample paragraph and every % comment with your own content.

% --- Introduction -----------------------------------------------------------
% State the problem, why it matters and what this document contributes.
\section{Introduction}
This document is a working example of the \textit{general} template. It shows
one use of each feature of this template, so that you can keep the parts you
need and replace the rest. The background material is summarised in a
textbook reference~\cite{example2}, and a short online note covers the same
topic~\cite{example3}. Replace this paragraph with your own introduction.

% --- Method -----------------------------------------------------------------
% Describe the model, the procedure or the algorithm that you follow.
\section{Method}
The sample method estimates a quantity from repeated measurements. Its
definition is the arithmetic mean of the observations,
\begin{equation}
  \bar{x} = \frac{1}{n} \sum_{i=1}^{n} x_i .
  \label{eq:mean}
\end{equation}
% Replace the equation above with the formula central to your work.
Every result in the next section is computed from this quantity.

% --- Results ----------------------------------------------------------------
% Report the main numbers first and explain what they mean afterwards.
\section{Results}
Table~\ref{tab:runs} lists three illustrative runs and the value and error
measured for each one.
% Replace the values, the caption and the label with your own measurements.
\begin{table}[H]
  \centering
  \caption{Illustrative runs, with the measured value and error}
  \label{tab:runs}
  \begin{tabular}{lcc}
    \toprule
    Run & Value & Error \\
    \midrule
    Baseline & 1.00 & 0.05 \\
    Variant A & 1.15 & 0.07 \\
    Variant B & \textbf{1.24} & 0.06 \\
    \bottomrule
  \end{tabular}
\end{table}

The direction of the calculation is often clearer in a drawing. The diagram
below is drawn from its source when the document is built, so the project
carries no image file.
% Replace the graph with your own; keep style= and the caption. The
% rankdir key is accepted by spell-whitelist.txt in the project root.
\begin{graphviz}[style=editorial, pos=H, width=0.62\linewidth, caption={From a measured run to its reported value}]
  digraph G {
    rankdir=LR
    run -> sample
    sample -> value
  }
\end{graphviz}

% --- Implementation ---------------------------------------------------------
% Show the few lines of code that a reader needs to reproduce the result.
\section{Implementation}
% Replace the listing with your own code. Keep it short (at most twelve
% lines) and set lang= to the language of the block.
\begin{code}[lang=python]
def mean(values):
    """Return the arithmetic mean of the values."""
    if not values:
        return 0.0
    return sum(values) / len(values)
\end{code}

% --- Conclusion -------------------------------------------------------------
% Restate the main result, name the limitations and the next step.
\section{Conclusion}
Replace this closing with a short summary: what the results show, which
limitations remain and what should be done next. The sample method follows
the framework described in~\cite{example}.
